\chapter[Anhang]{Anhang: Grafiken aus dem Versuch.}
\label{ch:Anhang}

Auf den Folgenden Seiten sind zu allen Messungen (außer der Ungekoppelten) Die Signale der Hallsonde und die FFT  Grafiken zu entnehmen. Insbesondee sind die Grafen zu den Signalen der Schwebung wichtig, da diese die in Inkscape bestimmten Größen beinhalten. 
Und ja, ich habe jeden Plot neu eingefärbt, damit diese in das Farbschema dieses Dokuments passen.


% --------------------------
% Symmetrische Schwingung
% --------------------------

\twocolumn
\onecolumn
\section{Symmetrische Schwinung}

\subsection[Schwache Kopplung]{Symmetrische Schwinung unter schwache Kopplung}

\img[1]{sym/Symmetrisch_S1-Weak_Signale}{Symmetrische Schwingung S1 bei schwacher Kopplung, aufgenommen mit der Hallsonde.}
\img[1]{sym/Symmetrisch_S2-Weak_Signale}{Symmetrische Schwingung S2 bei schwacher Kopplung, aufgenommen mit der Hallsonde.}
\img[0.5]{sym/FFT-batch_Page_7}{FFT der symmetrischen Schwingung bei schwacher Kopplung.}

\twocolumn
\onecolumn
\subsection[Mittlere Kopplung]{Symmetrische Schwinung unter mittlerer Kopplung}

\img[1]{sym/Symmetrisch_S1-Mittel_Signale}{Symmetrische Schwingung S1 bei mittlerer Kopplung, aufgenommen mit der Hallsonde.}
\img[1]{sym/Symmetrisch_S2-Mittel_Signale}{Symmetrische Schwingung S2 bei mittlerer Kopplung, aufgenommen mit der Hallsonde.}
\img[0.5]{sym/FFT-batch_Page_1}{FFT der symmetrischen Schwingung bei mittlerer Kopplung.}

\twocolumn
\onecolumn
\subsection[Starke Kopplung]{Symmetrische Schwinung unter starker Kopplung}

\img[1]{sym/Symmetrisch_S1-Strong_Signale}{Symmetrische Schwingung S1 bei starker Kopplung, aufgenommen mit der Hallsonde.}
\img[1]{sym/Symmetrisch_S2-Strong_Signale}{Symmetrische Schwingung S2 bei starker Kopplung, aufgenommen mit der Hallsonde.}
\img[0.5]{sym/FFT-batch_Page_4}{FFT der symmetrischen Schwingung bei starker Kopplung.}

% --------------------------
% Asymmetrische Schwingung
% --------------------------
\twocolumn
\onecolumn

\section{Asymmetrische Schwinung}

\subsection[Schwache Kopplung]{Asymmetrische Schwinung unter schwache Kopplung}

\img[1]{Asym/Asym_S1_Weak}{Asymmetrische Schwingung S1 bei schwacher Kopplung, aufgenommen mit der Hallsonde.}
\img[1]{Asym/Asym_S2_Weak}{Asymmetrische Schwingung S2 bei schwacher Kopplung, aufgenommen mit der Hallsonde.}
\img[0.7]{Asym/FFT_Page_2}{FFT der asymmetrischen Schwingung bei schwacher Kopplung.}

\twocolumn
\onecolumn
\subsection[Mittlere Kopplung]{Asymmetrische Schwinung unter mittlerer Kopplung}

\img[1]{Asym/Asym_S1_Middle}{Asymmetrische Schwingung S1 bei mittlerer Kopplung, aufgenommen mit der Hallsonde.}
\img[1]{Asym/Asym_S2_Middle}{Asymmetrische Schwingung S2 bei mittlerer Kopplung, aufgenommen mit der Hallsonde.}
\img[0.7]{Asym/FFT_Page_3}{FFT der asymmetrischen Schwingung bei mittlerer Kopplung.}

\twocolumn
\onecolumn
\subsection[Starke Kopplung]{Asymmetrische Schwinung unter starker Kopplung}

\img[1]{Asym/Asym_S1_Strong}{Asymmetrische Schwingung S1 bei starker Kopplung, aufgenommen mit der Hallsonde.}
\img[1]{Asym/Asym_S2_Strong}{Asymmetrische Schwingung S2 bei starker Kopplung, aufgenommen mit der Hallsonde.}
\img[0.7]{Asym/FFT_Page_4}{FFT der asymmetrischen Schwingung bei starker Kopplung.}

% --------------------------
% Schwebung
% --------------------------
\twocolumn
\onecolumn

\section{Schwebung}

\subsection[Schwache Kopplung]{Schwebung unter schwache Kopplung}

\img[1]{schweb/S1-Weak_Page 8_Coupled}{Schwebungssignal S1 bei schwacher Kopplung; gemessen wurden Schwebungsfrequenz (rot/grün) und Pendelfrequenz (blau); pink zeigt den Maßstab.}
\img[1]{schweb/S2-Weak_Page 9_Coupled}{Schwebungssignal S2 bei schwacher Kopplung; gemessen wurden Schwebungsfrequenz (rot/grün) und Pendelfrequenz (blau); pink zeigt den Maßstab.}
\img[0.7]{schweb/FFT-Weak_Page 7_Coupled}{FFT der Schwebung bei schwacher Kopplung.}

\twocolumn
\onecolumn
\subsection[Mittlere Kopplung]{Schwebung unter mittlerer Kopplung}

\img[1]{schweb/S1-Middle_Page 1_Coupled}{Schwebungssignal S1 bei mittlerer Kopplung; gemessen wurden Schwebungsfrequenz (rot/grün) und Pendelfrequenz (blau); pink zeigt den Maßstab.}
\img[1]{schweb/S2-Middle_Page 2_Coupled}{Schwebungssignal S2 bei mittlerer Kopplung; gemessen wurden Schwebungsfrequenz (rot/grün) und Pendelfrequenz (blau); pink zeigt den Maßstab.}
\img[0.7]{schweb/FFT-Middle_Page 3_Coupled}{FFT der Schwebung bei mittlerer Kopplung.}

\twocolumn
\onecolumn
\subsection[Starke Kopplung]{Schwebung unter starker Kopplung}

\img[1]{schweb/S1-Strong_Page 4_Coupled}{Schwebungssignal S1 bei starker Kopplung; gemessen wurden Schwebungsfrequenz (rot/grün) und Pendelfrequenz (blau); pink zeigt den Maßstab.}
\img[1]{schweb/S2-Stong_Page 5_Coupling}{Schwebungssignal S2 bei starker Kopplung; gemessen wurden Schwebungsfrequenz (rot/grün) und Pendelfrequenz (blau); pink zeigt den Maßstab.}
\img[0.6]{schweb/FFT-Strong_Page 6_Coupled}{FFT der Schwebung bei starker Kopplung.}
